% Brief: very brief; academic papers only (no websites); say what we borrow or adapt.
% Keys are in refs.bib. Add a paper there before citing it.
\section{Related Work}
\label{sec:literature}

\paragraph{Coded exposure.}
Raskar et al.\ opened and closed a camera's shutter in a binary code during one exposure, turning
motion blur into an invertible operation~\cite{raskar2006coded}. Agrawal and Xu then chose codes
for both invertibility and for estimating the blur from the image~\cite{agrawal2009coded}. We borrow
the principle and the 52-bit code of~\cite{raskar2006coded}, but code the light instead of the
shutter.

\paragraph{Other ways to capture motion.}
Motion-invariant photography moves the camera so that the blur is the same for every
speed~\cite{levin2008motion}, and velocity-dependent codes adapt the code to the expected
speed~\cite{mccloskey2010velocity}. Both need a moving camera or prior knowledge of the speed; our
rig needs neither, and estimates the speed afterwards.

\paragraph{Rolling shutter.}
Reading a sensor row by row distorts moving objects~\cite{liang2008rolling}, and coded rolling
shutter photography exploits the row timing deliberately~\cite{gu2010coded}. We instead model each
row's own exposure window in the blur operator.

\paragraph{Reconstruction.}
We use Tikhonov-regularised least squares as a transparent baseline, and plug-and-play restoration
with a learned denoiser (DPIR)~\cite{zhang2022plug} as the main solver.
